%%% Hlavní soubor. Zde se definují základní parametry a odkazuje se na ostatní části. %%%

%% Verze pro jednostranný tisk:
% Okraje: levý 40mm, pravý 25mm, horní a dolní 25mm
% (ale pozor, LaTeX si sám přidává 1in)
\documentclass[a4paper,oneside,12p]{report}
%\documentclass[a4paper,oneside,12p,twoside]{report}
\setlength\textwidth{155mm}
\setlength\textheight{247mm}
%\setlength\oddsidemargin{00mm}
\setlength\topmargin{-10mm}
\setlength\headheight{0mm}
\setlength\oddsidemargin{05mm}
\setlength\evensidemargin{05mm}
\let\openright=\clearpage


%% Vytváříme PDF/A-2u
\usepackage[a-2u]{pdfx}

%% Přepneme na českou sazbu a fonty Latin Modern
\usepackage[czech]{babel}
\usepackage{lmodern}
\usepackage[T1]{fontenc}
\usepackage{textcomp}

%% Použité kódování znaků: obvykle latin2, cp1250 nebo utf8:
\usepackage[utf8]{inputenc}

%%% Další užitečné balíčky (jsou součástí běžnýsrov.stribucí LaTeXu)
\usepackage{amsmath}        % rozšíření pro sazbu matematiky
\usepackage{amsfonts}       % matematické fonty
\usepackage{amsthm}         % sazba vět, definic apod.
\usepackage{bbding}         % balíček s nejrůznějšími symboly
			   										% (čtverečky, hvězdičky, tužtičky, nůžtičky, ...)
\usepackage{bm}             % tučné symboly (příkaz \bm)
\usepackage{graphicx}       % vkládání obrázků
\usepackage{fancyhdr}				% možnost slylizovat záhlaví
\usepackage{fancyvrb}       % vylepšené prostředí pro strojové písmo
\usepackage{indentfirst}    % zavede odsazení 1. odstavce kapitoly
\usepackage[nottoc]{tocbibind} % zajistí přidání seznamu literatury,
\usepackage{icomma}         % inteligetní čárka v matematickém módu
\usepackage{dcolumn}        % lepší zarovnání sloupců v tabulkách
\usepackage{booktabs}       % lepší vodorovné linky v tabulkách
\usepackage{paralist}       % lepší enumerate a itemize
\usepackage{caption}				%	popisky
\usepackage{dirtree}				% strom souborů
\usepackage[bottom]{footmisc}  % poznámky pod čarou vespod
\usepackage{bibentry}
\nobibliography*

\usepackage{color}
\usepackage{forest}					% na rodokmeny
\usepackage{natbib}
%\setbibentrystyle{author}

\definecolor{pblue}{rgb}{0.13,0.13,1}
\definecolor{pgreen}{rgb}{0,0.5,0}
\definecolor{pred}{rgb}{0.9,0,0}
\definecolor{pgrey}{rgb}{0.46,0.45,0.48}

%%% Údaje o práci

\def\NazevSkoly{Gymnázium, Praha 6, Arabská 14}
% Název oboru včetně počátečního 'Obor'.
\def\NazevOboru{Dějepis}

% Název práce v jazyce práce (přesně podle zadání)
\def\NazevPrace{Židovké město v Třebíči}

% Název práce v angličtině
\def\NazevPraceEN{Jewish town in Třebíč}

% Název práce v němčině
\def\NazevPraceDE{Jüdische Stadt in Třebíč}

% Jméno auto
\def\AutorPrace{Havránek Kryštof 2.E}

% Rok odevzdání
\def\RokOdevzdani{2020}
% Měsíc odevzdání
\def\MesicOdevzdani{Květen}

% Vedoucí práce: Jméno a příjmení s~tituly
\def\Vedouci{PhDr. Lenka Dvořáková }

% Nepovinné poděkování (vedoucímu práce, konzultantovi, tomu, kdo
% zapůjčil software, literaturu apod.)
\def\Podekovani{%
\textbf{Poděkování}

TODO

PhDr. Lenka Dvořáková

Milan Krčmář (zjisit si titul!)
}

% Abstrakt (doporučený rozsah cca 80-200 slov; nejedná se o zadání práce)
\def\Abstrakt{%
TODO
}
\def\AbstraktEN{%
TODO
}
\def\AbstraktDE{%
TODO
}

% 3 až 5 klíčových slov (doporučeno), každé uzavřeno ve složených závorkách
\def\KlicovaSlova{%
{klíčová} {slova}
}
\def\KlicovaSlovaEN{%
{key} {words}
}


%% Balíček hyperref, kterým jdou vyrábět klikací odkazy v PDF,
%% ale hlavně ho používáme k uložení metadat do PDF (včetně obsahu).
%% Většinu nastavítek přednastaví balíček pdfx.
\hypersetup{unicode}
\hypersetup{breaklinks=true}

%% Definice různých užitečných maker (viz popis uvnitř souboru)
%%% Tento soubor obsahuje definice různých užitečných maker a prostředí %%%
%%% Další makra připisujte sem, ať nepřekáží v ostatních souborech.     %%%

%%% Drobné úpravy stylu

% Tato makra přesvědčují mírně ošklivým trikem LaTeX, aby hlavičky kapitol
% sázel příčetněji a nevynechával nad nimi spoustu místa. Směle ignorujte.
\makeatletter
\def\@makechapterhead#1{
  {\parindent \z@ \raggedright \normalfont
   \Huge\bfseries \thechapter. #1
   \par\nobreak
   \vskip 20\p@
}}
\def\@makeschapterhead#1{
  {\parindent \z@ \raggedright \normalfont
   \Huge\bfseries #1
   \par\nobreak
   \vskip 20\p@
}}
\makeatother

% Toto makro definuje kapitolu, která není očíslovaná, ale je uvedena v obsahu.
\def\chapwithtoc#1{
\chapter*{#1}
\addcontentsline{toc}{chapter}{#1}
}

% Trochu volnější nastavení dělení slov, než je default.
\lefthyphenmin=2
\righthyphenmin=2

% Zapne černé "slimáky" na koncích řádků, které přetekly, abychom si
% jich lépe všimli.
\overfullrule=1mm

%%% Makra pro definice, věty, tvrzení, příklady, ... (vyžaduje baliček amsthm)

\theoremstyle{plain}
\newtheorem{veta}{Věta}
\newtheorem{lemma}[veta]{Lemma}
\newtheorem{tvrz}[veta]{Tvrzení}

\theoremstyle{plain}
\newtheorem{definice}{Definice}

\theoremstyle{remark}
\newtheorem*{dusl}{Důsledek}
\newtheorem*{pozn}{Poznámka}
\newtheorem*{prikl}{Příklad}

%%% Prostředí pro důkazy

\newenvironment{dukaz}{
  \par\medskip\noindent
  \textit{Důkaz}.
}{
\newline
\rightline{$\square$}  % nebo \SquareCastShadowBottomRight z balíčku bbding
}

%%% Prostředí pro sazbu kódu, případně vstupu/výstupu počítačových
%%% programů. (Vyžaduje balíček fancyvrb -- fancy verbatim.)

\DefineVerbatimEnvironment{code}{Verbatim}{fontsize=\small, frame=single}

%%% Prostor reálných, resp. přirozených čísel
\newcommand{\R}{\mathbb{R}}
\newcommand{\N}{\mathbb{N}}

%%% Užitečné operátory pro statistiku a pravděpodobnost
\DeclareMathOperator{\pr}{\textsf{P}}
\DeclareMathOperator{\E}{\textsf{E}\,}
\DeclareMathOperator{\var}{\textrm{var}}
\DeclareMathOperator{\sd}{\textrm{sd}}

%%% Příkaz pro transpozici vektoru/matice
\newcommand{\T}[1]{#1^\top}

%%% Vychytávky pro matematiku
\newcommand{\goto}{\rightarrow}
\newcommand{\gotop}{\stackrel{P}{\longrightarrow}}
\newcommand{\maon}[1]{o(n^{#1})}
\newcommand{\abs}[1]{\left|{#1}\right|}
\newcommand{\dint}{\int_0^\tau\!\!\int_0^\tau}
\newcommand{\isqr}[1]{\frac{1}{\sqrt{#1}}}

%%% Vychytávky pro tabulky
\newcommand{\pulrad}[1]{\raisebox{1.5ex}[0pt]{#1}}
\newcommand{\mc}[1]{\multicolumn{1}{c}{#1}}


%% Titulní strana a různé povinné informační strany
\fancypagestyle{plain}{
\fancyhead[C]{}
\fancyhead[L]{Ročníková práce - Gymnázium, Praha 6, Arabská 14}
\fancyhead[R]{\textbf{Židovské město v Třebíči}}
\fancyfoot[L]{Vypracoval: Havránek Kryštof 2.E (Programování)}
\fancyfoot[C]{}
\fancyfoot[R]{\thepage}
\renewcommand{\headrulewidth}{0.4pt}
\renewcommand{\footrulewidth}{0.4pt}
}



\begin{document}



%%% Titulní strana práce

\pagestyle{empty}
\hypersetup{pageanchor=false}
%%\setlength\oddsidemargin{15mm}

\begin{center}

{\LARGE\bfseries\NazevSkoly}

\vspace{-18mm}
\vfill

{\LARGE\NazevOboru}

\vfill

\centerline{\mbox{\includegraphics[height=4cm]{../img/logo.png}}}

\vspace{-8mm}
\vfill

{\bf\Large ROČNÍKOVÁ PRÁCE}

\vfill


\vspace{15mm}

{\LARGE\bfseries\NazevPrace}


\vfill


Vypracoval: \hfill {\AutorPrace}

Vedoucí práce: \hfill {\Vedouci}

\vspace{15mm}
\MesicOdevzdani \ \RokOdevzdani

\end{center}



\newpage


\openright
%%\setlength\oddsidemargin{00mm}

\hypersetup{pageanchor=true}
\pagestyle{plain}
\pagenumbering{roman}
\vglue 0pt plus 1fill

\medskip\noindent
Prohlašuji, že jsem jediným autorem této práce, všechny citace jsou řádně označené a všechna použitá literatura a další zdroje jsou v práci uvedené.
Tímto dle zákona 121/2000 Sb. (tzv. Autorský zákon) ve znění pozdějších předpisů uděluji bezúplatně škole Gymnázium, Praha 6, Arabská 14 oprávnění k výkonu práva na rozmnožování díla (§ 13) a práva na sdělování díla veřejnosti (§ 18) na dobu časově neomezenou a bez omezení územního rozsahu.

\vspace{10mm}

\hbox{\hbox to 0.5\hsize{%
V ........ dne ............
\hss}\hbox to 0.5\hsize{%
Podpis autora
\hss}}

\vspace{20mm}

\newpage
%
% %%% Poděkování
%
\openright
%
\noindent
\Podekovani

\newpage


\openright

\vbox to 0.20\vsize{
\setlength\parindent{0mm}
\setlength\parskip{5mm}

Název práce:
\NazevPrace

Autor:
\AutorPrace

% Vedoucí práce:
% \Vedouci, \KatedraVedouciho

Abstrakt:
\Abstrakt

% Klíčová slova:
% \KlicovaSlova

\vss}\nobreak\vbox to 0.20\vsize{
\setlength\parindent{0mm}
\setlength\parskip{5mm}

% Opakování v angličtině.

Title:
\NazevPraceEN

Authors:
\AutorPrace

Abstract:
\AbstraktEN

\vss}\nobreak\vbox to 0.20\vsize{
\setlength\parindent{0mm}
\setlength\parskip{5mm}

% Opakování v němčině.

Titlel:
\NazevPraceDE

Autoren:
\AutorPrace

Abstrakt:
\AbstraktDE

\vss}

\newpage

\openright



\tableofcontents
\newpage


\chapter*{Úvod}
\addcontentsline{toc}{chapter}{Úvod}

TODO


\pagenumbering{arabic}
\setcounter{page}{1}

\chapter{Počátky města Třebíče}

\section{Charakteristika města Třebíče}

Dnes je Třebíč nepříliš velké město v jihovýchodní části kraje Vysočina.
Hrálo však důležitou roli v historii okolí.
O tom svědčí fakt, že má dokonsrov.ě památky zapsané v seznamu UNESCO, a to Třebíčskou židovskou čtvrť a baziliku svatého Prokopa.
Počátky jeho historie jsou spjaty s benediktýnským klášterem, dnešním zámkem.
Je možné, že na území Třebečského lesa\footnote{území, na kterém byla založena Třebíč} se nacházelo i prehistorické osídlení, o tom však nemáme důkazy\footnote{srov. \bibentry{Uhlir1978}: 14}.

\section{Vznik kláštera}

Dějiny kláštera se začínají psát koncem 11. století.
Jednalo se o dobu, kdy na českém i moravském území docházelo k řadě konfliktů mezi jednotlivými Přemyslovci a moravskými knížeti. Právě ty vedly následek založení kláštera.
Břetislav II. jako svého následníka jmenoval svého nevlastního syna Bořivoje II., místo brněnského knížete Oldřicha\footnote{ten měl být králem dle rodové tradice Přemyslovců}.
Kvůli této a dalším neshodám uspořádal roku 1099 král Břetislav II. tažení na brněnský hrad, což donutilo Oldřicha a jeho bratra Litolda, jenž spravoval znojemský úděl, k útěku.
Po dvou letech vlády Bořivoje se však Oldřich a Litold opět vrátili na své úděly, využili šance kdy se po smrti Břetislava II. Bořivoj II. musel vrátit do Prahy.

Již při cestě zpět si uvědomili, že západní hranice jejich údělu jsou příliš daleko od hradů v Brně a Znojmě.
Z tohoto důvodu v roce 1101 založili na území Třebečského lesa téměř na hranicích s Přibyslavickem benediktýnský klášter.
Klášter měl sloužit jako předsunuté opevnění a kolonizační centrum.
Již v roce 1101\footnote{srov. \cite{Uhlir1978}: 20} byl postaven kostelík sv. Benedikta, který o tři roky později vysvětil pražský biskup Heřman.
Dále roku 1109\footnote{srov. \cite{Uhlir1978}: 20} vysvětil biskup olomouckým Jan II. samotný klášterní chrám Nanebevztí Pany Marie.
Klášter ihned po založení dostával různá územní darování, a to jak od Oldřicha a Litolda, tak i od jejich nástupců.
Samotný Litold (\CrossOpenShadow 1112) a Oldřich (\CrossOpenShadow 1113) jsou dle slov Tomáše Pěšiny\footnote{Tomáš Pěšiny z Čechorodu -- biskup, dějepisec, autor díla Mars Moravicus} pochováni v klášteře.

\begin{figure}[h]
	\centering
	\begin{forest}
		for tree={
    	child anchor=west,
    	parent anchor=east,
    	grow=east,
    	draw,
    	anchor=west,
    	edge path={
     		\noexpand\path[\forestoption{edge}]
       		(.child anchor) -| +(-5pt,0) -- +(-5pt,0) |-
       		(!u.parent anchor)\forestoption{edge label};
    		},
  	}
  	[Břetislav I.
    	[Vrastislav II.
		 		[Břetislav II.
						[Břetislav \textbf{D\footnotemark}
						]
						[Bořivoj II. (nevlastní syn)
						]
					]
				]
			[Kondrád I. Brněnský
				[Oldřich Brněnský \textbf{D}
					[Vrastislav \textbf{D}
					]
					[Vladislav
					]
				]
				[Litold Znojemský \textbf{D}
					[Kondrád II. Znojemksý \textbf{D}
					]
				]
			]
		]
		\end{forest}
		\caption[Rodokmen panovníků, vlastní tvorba]{
		\centering
			Rodokmen panovníků figurujících v rané historii kláštera, písmeno \textbf{D} značí knížata, která darovala území klášteru
		}
\end{figure}
\footnotetext{Dle opisu Kroniky Kosmovy, daruje vsi Břetislav, nejspíše se jedná o tohoto}

Bohužel je jen málo pramenů z této doby, jako jediný zdroj o rané historii kláštera máme Třebíčský (Brněnský) opis Kosmovy Kroniky\footnote{srov. \cite{Uhlir1978}: 18}.
Do něho byla klášterním mnichem připsaná k roku 1115 zmínka o založení kláštera.
Mnich také uvádí jaké vsi byly postupně darovány klášteru.
Údaje jsou však sporné, a pravděpodobně vycházel z pozdějších zdrojů, neboť počet darovaných vsí je dle vpisku sedmdesát tři, což značně převyšuje i pozdější čísla.
Ve vpisku jsou uvedeny i vsi, u kterých je pozdější připojení písemně doloženo.
Existovala také klášterní kronika z roku 1201, jejímž autorem byl opat Tiburcius, ta se však ztratila.

\section{Tržní osady okolo kláštera}

Krátce po založení kláštera se začaly v jeho okolí vytvářet tržní osady.
Z nich se později zformovalo samotné město Třebíč.

\subsection{Podklášteří}

Vhodná poloha kláštera, v jeho blízkosti se nacházela řada cest, a také existence hradeb, které mohly ochránit poddané v době nepokojů, přispěly na přelomu 12. a 13. století k vzniku tržní osady.
Ta se samovolně začínala formovat přímo pod klášterem na levém břehu řeky, klášterní hradby nebyly totiž dost rozsáhlé, aby celou osadu pojaly.
Nově vzniklá osada Podklášteří ihned začala přitahovat kupce a nejspíše i Židy z širokého okolí, což vedlo k jejímu rychlému růstu.
Za nedlouho tak začala být omezována svou polohou -- skály bránily v jejím růstu a ohrožovaly ji záplavy řeky Jihlavy.
Podklášteří postupně ztrácelo na svém významu jakožto tržní osada.\footnote{srov. \cite{Uhlir1978}: 22}.

\subsection{Vznik Starého města}
%% přepsat
V 1. třetině 12. století vznikla na příkaz klášterní vrchnosti nová tržní osada, dnešní Staré město.
Ta vznikla na pravém břehu řeky Jihlavy, kde bylo více prostoru pro rozvoj, než v již zastaveném Podklášteří.
Osada ihned začala nabývat na významu a tvořila se okolo ní aglomerasrov.lších osad, můžeme zmínit ves na Polánce, či vsi Kosovice, Arklebice a Horku, jenž opat Lukáš roku 1225 vyměnil se znojemskou šlechtičnou za ves Oslavany, které se nacházely dál od kláštera.
Stejně jako Podklášteří byla i nová osada omezena svým umístěním, což s postupem času opět začalo brzdit vývoj na území.\footnote{srov. \cite{Uhlir1978}: 26}
Proto ještě v 1. polovině 13. století bylo založeno dnešní Nové město.


\section{Založení města Třebíč}
Město Třebíč bylo podobně jako Staré město založeno z vůle opata, jenž si vyžádal povolení u krále Přemysla Otakara II.
Z toho důvodu se také o městu Třebíč hovoříme v této době jakožto o poddanském městu náležící církevní vrchnosti.
Opat dále musel vést vyšší soudy, či rozhodoval o manství\footnote{nápravníci, lidé disponující svobodným lénem}.
Opět se jednalo o strategický krok, města totiž disponovala různými privilegii, jako právo trhu či právo hradební, právem vářečným již disponovala Staržečska\footnote{Jiné označení Starého města, dnes čtvrť v Třebíči}.
První písemná zmínka o městu Třebíč pochází z roku 1277, v ní opat Martin oceňuje lokátora Heřmana\footnote{srov. \cite{Uhlir1978}: 27}, který město Třebíč založil.

\subsubsection{Přestavba kláštera}

Během 13. století dochází k rozsáhle přestavbě kláštera.
Byly vztyčeny definitivní klášterní budovy a roku 1260 byla dokončena stavba klášterního chrámu.
Společně s přestavbou budov bylo přestavěn a vylepšen i hradební systém okolo kláštera.
Bohužel během Bitvy o Třebíč\footnote{12. květen — 9. červen 1468, součást česko-uherských válek (Druhá válka husitská)} byla většina staveb zničena či poškozena.
Postupem času vznikaly i další budovy patřící klášteru rozeseté po celém klášterem spravovaném území.

\chapter{Pozdější vývoj města Třebíče}

V následující kapitole se zaměříme na pozdější vývoj města Třebíče.
A to od bitvy na Moravském poli až po konečné vítězství českých stran nad německými během národního obrození.
Třebíč se poté již vyvíjí relativně stabilně stejně jako ostatní moravská či česká města.
Dění obou válek se jí až na lidské (vyhoštění židů) a finanční ztráty příliš netýká.
Stejně tak vzrůst komunismu a jeho následný pád ji postihne stejně jako jiná česká města.

\section{Bitva na moravském poli až nástup Pernštejnů}

Klid v kraji neměl dlouhého trvání.
Kvůli bitvě na Moravském poli (26.srpen 1278) vzrostla moc šlechty, která pořádala loupeživé výpravy s cílem zmocnit se královského majetku, ohrožena byla i samotná Třebíč.
Situaci zachránil až nový král Jan Lucemburský, který sám Třebíč navštívil.
Město poté dále prosperovalo a jeho význam pomalu zastínil i samotný klášter.
Výrazně pozdějšímu vývoji pomohlo privilegium markraběte Karla, budoucího krále Krala IV, díky kterému město získalo pevný právní řád.\footnote{srov. \cite{Uhlir1978}: 38}

Je nutno zmínit že v této době hovoříme o Moravě jako o markrabství.
Tím se stala díky Fridtichu I. Barabarosovi, který na směnu v Řezně markrabství deklaroval.
Uskutečnitlo se tak v rámci řešení sporu mezi knížetem Bedřichem a Konrád III. Ota Znojemský, pozdější kníže Konrád II. Ota.
Kdy právě Ota se stává roku 1182 prvním markrabětem.
Je zajímavé, že markrabský titul vyjadřoval podřízenost českému králi, nikoliv králi Svaté říše římské.\footnote{srov. \bibentry{Wihoda2010}: 237, 241, 242}

Do dalších problémů se dostává město po smrti markraběte Jana Jindřicha \CrossOpenShadow 1378, té následovaly konflikty mezi Joštem a Prokopem, kteří vzájemně soupeřili o kontrolu nad Moravou.
Tyto spory vyústily až v domácí válku, která koncem 14. století vrcholí těžkou hospodářskou krizí.
Situaci uklidnil až král Václav IV. roku 1413.
Doba klidu však neměla dlouhého trvání, protože již o dva roky později se začal projevovat odpor proti církvi jako ohlas na kázání mistra Jana Husa. \footnote{srov. \cite{Uhlir1978}: 28--43}

O samotném vývoji husitství v Třebíči nemáme žádné prameny, je však známo, že Třebíč udržovala vztahy se moravskými šlechtici, jež byli myšlenkám Jana Husa nakloněni.
Během Žižkova tažení na Moravu se třebíčský klášter dostal pod kontrolu husitů, z kláštera poté podnikaly výboje dál do okolí.
Konečně se situace zklidnila až v roce 1435 díky bitvě u Lipan, k které došlo v předchozím roce.
Během hustiských válek došlo k výraznému hospodářskému úpadku, ze kterého se klášter dostával celá desetiletí.
Třebíč byla po válkách věrný králi Jiřímu z Poděbrad. \footnote{srov. \bibentry{Uhlir1978}: 49--54}

Již 12. května 1468 se Třebíč dostala do obležení, a o dva dny později byla Matyášovými (uherský král Matyáš Korvín) vojsky dobyta.
Velká část obyvatel (okolo čtyř tisíc lidé) se však uchýlila do opevněného kláštera.
6. června se česká vojska pokusila klášter zachránit, ale bez úspěchu, 15. června se posádka v klášteře definitivně vzdala.
Po celých sedm let bylo město opuštěné. \footnote{srov. \bibentry{Uhlir1978}: 54--46}

\section{Vláda Pernštejnů}

S obnovou města se začalo teprve v roce 1474 přímo na troskách původní Třebíče.
Celková restaurace města a klášterního velkostatku však nebyla možná kvůli obrovským dluhům, kterým klášter čelil.
Nejvíce obnově napomohl Vilém z Pernštejna\footnote{bohatý rod, ve službách Jagellovců, třebíčský statek byl pod jeho kontrolou}, ten také v kraji zavedl rybnikářství a započal opravy samotného kláštera.
Roku 1507 Vilém rozdělil své území mezi syny. V závěti také apeloval na nutnost oprav různých klášterů v okolí o třebíčském nepadla ani zmínka.
Přítomnost benediktýnů mu bránila v převedení Třebíče do svého dědičného vlastnictví.
Dále se Třebíče ujímá Jan IV. z Perštejna, který pokračuje v otcově politice.
Na krátkou však dobu město připadá pod vládu Jana Jetřicha\footnote{Zeť Jana z Pernštejna, menší šlechtic, Třebíč byla jeho největším panstvím. Vede část české jednotky v tureckých válkách}, který jako první světský správce razantně zasahuje do vedení města.
Jan z Perštejna přiděluje městu další práva, z nichž je nejdůležitější právo mílové, které značně zlepší ekonomické postavení města.
Obnovuje také na žádost městské rady dekret Jana Jetřicha Černohorského, který omezoval konkurenceschopnost židovské populace ve městě.
Po Janu (\CrossOpenShadow 1548) se vlády ujímá syn Vratislav.
Ten již Třebíč převádí do dědičného vlastnictví. \footnote{srov. \bibentry{Uhlir1978}: 61--71}

\section{Vláda Osovců}

Mezi polovinou 16. století a rokem 1613 byl Třebíč državou Osovců, poté co Burian Osovský odkoupil třebíčské panství.
Vzhledem k ceně třebíčského panství se však Burian zadlužil, to se snažil vykompenzovat produkcí piva.
Omezoval také různá práva ve městě jako je -- právo na pronájem, ale i právo mílové.
Město proto podalo stížnost na zemský soud, ta však neuspěla, Třebíč se tedy odvolala k císaři.
Ferdinant I. poté 7. března 1560 obnovil platnost všech privilégií, zároveň napomenul Buriana, ten již dále nepodnikal tak razantní kroky proti městu. \footnote{srov. \cite{Uhlir1978}: 77--78}

Roku 1568 se vlády nad městem ujal Smil Osovský z Doubravice\footnote{Jeden z důležitých pramenů doby, své paměti si zapisoval do Codexu Dubraviciana}.
Smil musel čelit dluhům svého otce, ty částečně vyřešil sňatky, ale hlavně díky razantnímu zlepšení ekonomické situace ve městě.
Zlepšení však bylo na úkor poddaných, a to v podobě vyšších rent, či delších robotních povinností.
Snažil se ale i postarat se o sirotky, předcházet požárům, zabezpečit učitele, podporovat také městský špitál, či dávat mírnější tresty -- například za pytlačení nebyla vydloubnuta oči.
Roku 1573 dále vydal Zřízení selské, které nabádalo k poslušnosti vůči úředníkům a zakazovalo sběr dříví a pastvu v panských lesích.
O devět let později získalo město právo pečetit červeným voskem a hned v příštím roce bylo vydáno městské zřízení, na jehož základě bylo město zpravováno pánovými instrukcemi.

Třebíč byla počátkem 17. století opět prosperujícím městem s rozvinutou řemeslnou výrobou a jejíž správce již nebyl zadlužen.
Během správy Smila Osovského také došlo k přestavbě kláštera na zámek.\footnote{srov. \cite{Uhlir1978}: 78--96}

Po úmrtí Smila se vlády ujímá vdova Kateřina z Valdštejna, ta se sice provdala za Karla staršího ze Žerotína, ale Třebíč si nechala pod svojí kontrolou.
K městu se chovala mírně a udržovala dobré vztahy s měšťany. \footnote{srov. \cite{Uhlir1978}: 101--102}


\section{Třicetiletá válka}

I přes nátlak ze strany krajanů zůstal Žerotín nestranný a defenestraci Čechů neschvaloval jako adekvátní.
Žerotín nepodpořil ani krále Fridricha, zimního krále, v jeho úspěch totiž nevěřil, naopak samotná Třebíč zimnímu králi nabídla nocleh při jeho cestě do Brna a předala dary.
O dalším předbělohorském vývoji nejsou žádné prameny.
Karel starší ze Žerotína také slíbil ochranu města a 9. 12. 1620 sám císař Ferdinant II. zaručil, že císařská armáda město nepoškodí a bude ho chránit.
Mezi 14.-15. prosincem 1620 tak městem císařská armáda jen prošla a zanechala zde ochrannou posádku.
Roku 1621 proto Třebíč dokázala vyhovět prosbám o pomoc měst Jihlavy a Znojma, která pod ochranou císaře nebyla. \footnote{srov. \cite{Uhlir1978}: 103--106}

Počátkem ledna 1622 do města přijel Marradasův jezdecký pluk císařské armády o síle pěti set jezdců, který představoval pro Třebíč značnou zátěž.
Pluk se navíc nijak nehrnul do bojů a v městě pobyl až do října.
Potenciální nebezpečí pro město představoval Gábor Bethlen poté co začal podnikat nájezdy na moravské území a pustošit místní města, Třebíč proto požádala o pomoc.
Gábor však začal mírová jednání předtím, než byla Třebíč přímo ohrožena.
Mezitím však dorazily dva regimenty slezského vojska o síle tří tisíc vojáky, ti v Třebíči tábořily dvacet tři dní.
Roku 1626 došlo k dalšímu uherskému výboji, což vyvrcholilo opět v pobyt císařských vojsk v Třebíči. \footnote{srov. \cite{Uhlir1978}: 106--107}

Další ránu městu zasadilo Obnovené zřízení zemské.
Po něm již náboženskou svobodu, kterou Třebíč dříve disponovala,  nemohla šlechta nijak garantovat\footnote{Na Moravě byla náboženská svoboda dlouho standardem, proto se jí netýkal ani Rudolfův majestát.}.
Brzy po vydání emigračního patentu (9. březen 1928) se proto Žerotín rozhodl pro emigraci, konkrétně do Vratislavi\footnote{Činní tak v rámci pomoci nově persekuovaným, zvláště Jednotě bratrské}.
Z Vratislavi mohl stále spravovat území, které si ponechal, a Kateřina stále byla schopna dohlížet na Třebíčí.
Novým držitelem Třebíče se na jaře roku 1928 stává Kateřin bratr Adam z Valdštejna, samotná Kateřina roku 1638 zemřela a o dva roky dříve zemřel Karel starší ze Žerotína. %
Jejich těla jsou pochována v Třebíči. \footnote{srov. \cite{Uhlir1978}: 108--110}

Adam však nejevil o správu města velký zájem a přenechal ji synu Rudolfovi z Valdštejna.
Ten již neměl takové postavení politické jako jeho otec a ani příliš nezasahoval do náboženských poměrů ve městě.
Otázky náboženství nepovažoval za příliš důležité a spíše se zajímal o otázky ekonomické, které řešil vysokými poplatky. \footnote{srov. \cite{Uhlir1978}: 109--112}

Vývoj třicetileté války během přechodu z dánské fáze do fáze švédské příliš Třebíč neovlivnil.
Až tažení na Moravu Linharta Torstensona\footnote{Švédský polní maršál, velmi schopný a obratný vojevůdce.} Třebíč postihlo
Roku 1644 se musela podílet na vydržování a vystrojení vojsk, která bránila Olomouc.
V březnu 1645 však do města přišla švédská jednotka, ta si vynutila 6500 zl. výpalného a obnovila v nekatolické bohoslužby, musela také platit poplatek švédské posádce v Jihlavě. \footnote{srov. \cite{Uhlir1978}: 112--113}

\section{Obobí rekatolizace}

Moravská šlechta pod vedením Jana Rottala\footnote{Moravský šlechtic z rodu Rottalů, držitel řádu zlatého rouna a jedna z významných osobností habsburské monarchie.} začala vydávat jedno nařízení proti nekatolíkům za druhým, stejně se choval i císař se svými pateny z roku 1651.
V téže roce byla i sestavena komise s cílem provádět inspekci církevního majetku na Moravě.
Stejně jako pro další moravské města tak i pro Třebíč znamenala druhá polovina 17. století období rekatolizace až do 18. století však tajné nekatolictví v městě vydrželo.
Došlo také k obnovení cechů, které postupně eliminovaly konkurenci a začali brzdit vývoj v městě. \footnote{srov. \cite{Uhlir1978}: 125--126}

Samotné město se dostalo pod správu Adama Františka, jeho ze začátku zastupovala Zdislava ze Sezimova Ústí.
O Zdislavině působení nemáme příliš písemností, víme však, že odpustila dluh kupcově rodině a vyhověla stížnosti Židů, aby nebyli vystavováni násilí nestihnou--li robotu.
30. června 1655 se vlády ujal syn Adam, ten se choval podobně jako jeho otec a otázce náboženství tak nepřikládal velkou hodnotu.
Již však nebyl schopen odolávat náporu z venčí a zvláště žalobě polenského děkana Kryštofa Kazimíra Burešovského, ta byla podložena řadou důkazů obsahovala však i polopravdy, či naprosté lži, jako bylo obvinění z existence 21 sekt v Třebíči.
Právě tato žaloba vyvolala na úřadech velkou reakci, což vyvrcholilo vytvořením \uv{inkviziční komice} zorganizované 2. prosince 1656 novoříšským proboštem Engelbertem a moravskobudějovickým Rudolfem Jindřichem ze Schaumburka.
Ta působila ve městě od 3. do 9. února 1657, za přítomnosti rychtáře, faráře a i samotného žalobce. \footnote{srov. \cite{Uhlir1978}: 126--128}

Komise například zjistila, že v Podklášteří žilo 13 nekatolických rodin ku 7 katolickým a podobně na tom byl zbytek Třebíče.
Nejvísrov.minovali utrakvisté, vyskytovali se však i luteráni.
Závěr komise je byl relativně jasný -- Třebíč je útočištěm nekatolíků, zvláště zběhlých z Čech.
Komise si také stěžovala na chování židů a obvinila také město z útlaku katolíků, proti tomu se však městská rada ihned ohradila.
Děkan Burešovský a farář Slovatius navíc přidali další stížnosti vůči Třebíči -- závěrečná zpráva tak pro Třebíč nebyla nikterak dobrá.  \footnote{srov. \cite{Uhlir1978}: 128--129}

Již 8. května 1657 přichází do města protireformační komise vedená hejtmanem Tannazollem.
Její prvotní snahy však byly neúčinné, celkově od listopadu 1658 do července 1659 se podařilo obrátit na \uv{pravou} víru jen 109 lidí, obecně lidé neměli přílišnou motivaci, pokuty pro nekatolíky nikdo nevymáhal.
Přelomem 50. a 60. let, i přes nezdárné úspěchy komise, došlo k nová vlně emigrace a 20\% šlechty opustilo Třebíč. \footnote{srov. \cite{Uhlir1978}: 130--131}

Počátkem 60. let začalo také hrozit nebezpečí se strany Turků, kteří pořádali vpády na Moravu.
Roku 1664 však porážkou Turků u Svatogotthardského průsmyku ohrožení pominulo. \footnote{srov. \cite{Uhlir1978}: 132--133}

Po Adamově smrti (\CrossOpenShadow 1666) připadá město Františkovi Augustinovi.
Za jeho vlády se dostává rekatolizasrov. nové fáze, František zároveň usiluje o zřízení kapucínského kláštera, jeho žádosti však nebylo vyhověno.
Jeho správa neměla dlouhého trvání a brzy se na jeho místo dostává bratr Karel Ferdinand.
Ten se, na rozdíl od jeho předchůdců, zapsal jakožto dobrotivý pán.
Během jeho vlády do města přicházejí i zástupci kapucínů a rychle se začleňují do městské společnosti. \footnote{srov. \cite{Uhlir1978}: 133--135}

\section{Spory mezi měšťany a vrchností}

Za vlády Josefa I. docházelo k řadě sporů mezi měšťany a vrchností.
Ty zakončil až reskript císaře z 15. prosince 1708.
V něm císař ponechal Třebíči soudní pravomoc, zároveň však apeloval, aby se Valdštejn na jurisdikci podílel a zaručil také měšťanům právo na svobodný trh. \footnote{srov. \bibentry{Janak1981}: 28--29}

I tak roky 1724 a 1726 došlo k dalším sporům mezi městem a aktuálním správcem -- Janem Josefem z Valdštejna.
Ty vyvrcholily v narovnání roku 1726, které posunulo město dál k emancipaci měšťanů. \footnote{srov. \cite{Janak1981}: 29--33}

K dalšímu sporu došlo roku 1728, kvůli termínu tělesné člověčenství \footnote{termín k vyjádření závislosti na feudálovi, stejné jako \uv{nevolnictví}}, jenž v jedné listině Jan Josef použil, tento spor měšťané prohrály \uv{nevolnický}.
Spor se dostal před tribunál ještě jednou díky kupci Janu Bernardovi Matesovi, který odmítl přijmout \uv{nevolnictví}, načež mu bylo odebráno právo obchodovat, to bylo v rozporu s právem na svobodný trh, vrchnost zde již byla nucena ustoupit.
Po smrti Jana Josefa se stal držitelem města František Heřman Arnošt z Valdštejna, ten se výrazu nevolnictví plně zřekl a nepoužily ho ani jeho následovníci. \footnote{srov. \cite{Janak1981}: 33--35}
Spory se vyskytovaly i za vlády Marie Terezie.

\section{Vláda Marie Terezie, Josefa II a napoleonské války}

Za vlády Marie Terezie se postupně začalo opouštět od tradiční ekonomice feudální společnosti a jít vstříc kapitalistickému fungování.
V Třebíči se vývoj odehrával podobně a na důležitosti nabyli zejména tkalci, ševci a koželuži, oproti tomu upadlo pro Třebíč dlouhou dobu charakteristické soukenictví, které dostalo do problémů z důsledku nedostupnosti vlny, což mělo za následek i soudní spor.
Můžeme zmínit, že obchodem se živila většina židovských obyvatel v roce 1773 to bylo 55, v roce 1799 již 92.
Samotná průmyslová revoluce však přišla do Třebíče trochu opožděně. \footnote{srov. \cite{Janak1981}: 13--29}

Třebíč se stala obětí válek o děditství rakouské, kdy 10. února 1742 do města vtrhli pruští husaři a muselo jim být zaplaceno 12 000 zl. výpalného a 10000 zl. navíc od židů.
Dále 12. února dorazila armáda krále Bedřicha Velikého o síle 12 000 mužů, kterou město muselo ubytovat a 24. února.
Devastující byly také neúrody mezi léty 1765 a 1769, dále 1771 a 1772, ty vyvrcholily až v selské povstání z roku 1775. \footnote{srov. \cite{Janak1981}: 38}

1. ledna 1755 přišla Třebíč o právo hrdelní.
Tato změna souvisela s reformou, která omezovala počet institucí, jenž tento trest mohly vynést. \footnote{srov. \cite{Janak1981}: 39}
Samotná Marie Terezie do vedení měst příliš nezasahovala, to změnil až její syn -- Josef II.

První Josefův zásah byl prostřednictvím nového dekretu o soudnictví z 27. prosince 1786, podle kterého mělo být soudnictví soustředěno do kriminální soudů, jako které fungovaly magistráty krajských měst, v případě Třebíče to byla Jihlava.
Také došlo k zrušení městských rad a jejich nahrazení magistráty, jejichž zastupitelé -- purkmistr a rada volilo měšťanstvo prostřednictvím měšťanského výboru, samotné kandidáty však registrovala vrchnost, což jí dávalo větší moc, než jakou mělo měšťanstvo.
Spory s vrchností však neustaly ani po vytvoření magistrátu (1. volby se konaly v březnu roku 1787), vrchnost také kromě schvalování kandidátů rozhodovala o výdajích města.
Docházelo také k silné germanizaci ve správě města, a tak veškerý spisy z této doby jsou v němčině. \footnote{srov. \cite{Janak1981}: 36--37}

Dne 17. srpna 1796 vydal František I. Rakouský\footnote{císař rakouský František II.} patent vyzývající občany k vojenské službě, samotné město dalo na válečnou snahu značný peněžní příspěvek.
To i přesto, že na přelomu 18. a 19. století bylo obecní hospodářství ve velmi špatném stavu.
Mezi 19. a 26. listopadem 1805 prošel městem první sbor pod vedením maršála Bernadottema, toho uvítal farář Dvořecký a nabídl Bernadottemu nocleh na zámku, ten se však spokojil s domem radního.
Za nedlouho však ve městě vypukl nedostatek jídla a francouzští vojáci museli táhnou dál.
Dále se ještě musela Třebíč postarat o dvě divize, které sem přišli po bitvě u Slavkova a pobývaly v Třebíči mezi 9. až 29. prosincem.
Celkové výdaje dosáhly 300 000 zl., ztráta nikdy nebyla vykompenzována.
Situaci navíc ještě zhoršily špatné úrody v letech 1807 a 1808. \footnote{srov. \cite{Janak1981}: 38--39}

Po uzavření příměří po bitvě u Znojma Třebíč protínala demarkační linie, což způsobilo městu další finanční ztráty o celkové hodnotě 100 000 zl.
Navíc do města byly svezeni ranění obou bojujících stran, to mělo za následek tyfovou epidemii. \footnote{srov. \cite{Janak1981}: 39}

Další devastující událostí pro Třebíč byl státní bankrot z roku 1811, ten obral malé živnostníky i o to málo co jim zbylo a navíc mezi vrchností se rychle rozšířila korupce.
Židé navíc byli osočeni, že o bankrotu věděli předem a začali předem skupovat majetek.
Situaci ještě zhoršily neúrody v letech 1815 a 1817 a vše završila morová epidemie mezi dobytkem roku 1819, po které přestala fungovat většina řemeslné výrovy ve městě. \footnote{srov. \cite{Janak1981}: 39--40}

Načež roku 1821 proto došlo k robotní vzpouře.
Vzbouřenci však byli odvedeni na zámek a zbiti, městem poté začaly kolovat letáky vyzývající k ukončení trestů s výhrůžkou žhářství.
Hrozby se staly skutečností a dne 3. května bylo vypáleno 100 domů v židovské obci a dalších 50 v Podklášteří. \footnote{srov. \cite{Janak1981}: 30}

V červnu roku 1822 město opět zachvátil požár, který měl více devastující účinky než bitva o Třebíč -- shořelo přes 290 domů a to včetně kostela, radnice, fary a jedné školy.
Třebíč se poté opět dostala do finančních problému, které vyřešil až syndik Gottsmann.

Třebíči však nebylo dovoleno dlouhému oddychu, v roce 1830 byla Třebíč zasažena povodní, v roce 1832 a opět v roce 1836 vypukla epidemie cholery a v roce 1847 město opět zasáhl požár.


\section{Industrializace a národní hnutí v Třebíči}

Národní hnutí se v Třebíči začalo projevovat ve 20. letech 19. století.
Většina obyvatel Třebíče byla česká, a případní cizinci byly do českého prostředí rychle asimilováni, správu nad městem však drželi němečtí úředníci, a tak se český živel prosazoval jen pozvolna.
Prvním vlastencem jehož chování mělo nějaký dopad se stal mladý lékař Jan Miloslav Haněl, který do Třebíče přišel v roce 1835.
Ten o devět let později založil čtenářský spolek a angažoval se také v dobrovolnické divadelní činnosti, kterou pořádal spolu s českými studenty. \footnote{srov. \cite{Janak1981}: 33--35}

Samotné události roku 1848 zastihli vedení Třebíče nepřipravené, situaci stabilizovalo až vytvoření národní gardy, to ale také vedlo k rozdělení měšťanstva na stranu německou, v jejímž čele stál Ignác Glas, a na stranu českou s dr. J. M. Hanělem v čele.
Zatím však národní garda pracovala společně s Němci. \footnote{srov. \cite{Janak1981}: 69--70}

1. českou manifestací se stalo svěcení praporu národní gardy uspořádané dne 28. května 1849, při něm se zpívaly české písně a provolávala různá hesla.
Garda však neměla dlouhého trvání a již koncem roku 1849 musela z nařízení ukončit svoji činnost. \footnote{srov. \cite{Janak1981}: 70--72}

20. října 1860 byl vydán Říjnový diplom, ve kterém byla zaslíbena konstituční vláda, načež 21. února 1861 byla vydána samotná ústava.
Obě události měšťané přijali s nadšením, a záhy se začala angažovat v Třebíči česká národní strana a vytvořila spolek Měšťanská beseda s dr. J. M. Hanělem v čele.
Jednalo se o jeden z prvních podobných národních spolků na území Moravy.
Německá Třebíč jako odpověď na český spolek utvořila organizace Gesangverein, Turnverein a Schützverein (původní střelecký spolek). \footnote{srov. \cite{Janak1981}: 76}

Období 40. let se také neslo v duchu sociální nejistoty drobných výrobců, ta vyvrcholila vystoupeními namířenými hlavně proti židovské čtvrti.
Během let 48-49 došlo k dočasnému zrovnoprávnění, což vedlo i k přesun zámožných židů z židovského města na náměstí.
To však mělo za následek další bouře proti židovskému obyvatelstvu, k těm došlo mezi 1. dubnem a 7. květem 1850.
K dalším nepokojům došlo v roce 1860, společně úpadkem pro Třebíč charakteristického soukenictví.
Město také trpělo akutním nedostatkem peněz, byly proto vydány dočasné bankovky, ty se však velmi lehce padělaly.
Celá řada měšťanů se dostala do finančních problémů a byla nucena prodat své domy.
To mělo za následek další proti židovské nepokoje (1863 a 1866). \footnote{srov. \cite{Janak1981}: 77--78}

\section{Prusko-rakouské války}

Během prusko-rakouské války došlo k poslednímu vzrůstu soukenictví díky armádním zakázkám,ty doprovázely i rekvizice, které město přišli až na 250 000 zl.
Třebíč se navíc dočkala opět nepřátelské invaze, a to 13. července, kdy přitáhla pruská vojska o síle 35 000 mužů.
Ta zde pobyla do 16. prosince.
Dále město hostilo záložní sbory a zdravotní kolony.
Okupace trvala až do uzavření příměří, poslední jednotky odešli až 8. září.\footnote{srov. \cite{Janak1981}: 78}

\section{Další vývoj}

20. srpna 1867 Třebíč dostala svého českého starostu -- Ignáce Přerovského, zároveň radními příslušníky se stali členi Měšťanské besedy jako byl J. M. Haněl.
Vznikla také Občanská záložna, která měla fungovat jakožto finanční opora českého národního hnutí.
K založení Sokola však v Třebíči nedošlo. \footnote{srov. \cite{Janak1981}: 79--83}

Při dalších volbách roku 1870 Německá strana vůbec nekandidovala a do obecního výboru tak byly zvoleni jen čeští funkcionáři.
Situace však neměla dlouhého trvání, již v roce 1871 se moravským místodržitelem stal baron Weber, který opět začal potlačovat národní snahu, a rozpustil i moravský zemský sněm.
Němců se zastala i většina židovského města a v dalších volbách byl jako starosta zvolen představitel německých liberálů -- dr. Jan Bažant.
Krátce po jeho zvolení byla definitivně ukončena činnost záložny, a o rok později bylo zrušeno i obecní zastupitelstvo, na základě obvinění z mrhání penězi. \footnote{srov. \cite{Janak1981}: 82--83}

Správu nad městem poté dočasně převzala komise a ta se postarala o uspořádání voleb.
Ty byly zmanipulované v neprospěch České národní strany.
7. prosince 1873 bylo ustanoveno nové obecní zastupitelstvo -- Česká národní strana skončila s menšinou a do samotné městské rady se nedostal jediný český zástupce.
Příštích voleb se už čeští zástupci nezúčastnily vůbec a ve volbách v roce 1879 bylo vedení města opět z většiny německé. \footnote{srov. \cite{Janak1981}: 83--90}

% byly změny v židovském městě a samotní židé hodně ovlivňovaly politiku Třebíče, budu se věnovat v samostatné kapitole o městě

Ke změně došlo až u voleb z roku 1879, kdy česká strana získala naprostou většinu hlasů a Třebíč opět dostal českého starostu.
Jednalo se o první velké město na Moravě, kde se Čechům něco podobného podařilo.
Třebíč však byla opět ve finančních problémech, německá správa zanechala dluh 107 727 zl. (celkový dluh činil 123 455 zl.).
Byla proto zvýšena přirážka u přímých daní na 12\%, a to k velké nelibosti pracujících.
Stát zároveň začal dávat větší důraz na místní samosprávu, věc na kterou město vůbec nemělo prostředky, přirážky se tak i přes nepokoje s musely zvyšovat\footnote{srov. \cite{Janak1981}: 93--94}

Konec 19. století a začátek 20. se vyznačoval tvorbou různých spolků a hnutí, mezi nejdůležitější patřilo dělnické, holdující sociálně demokratické straně a operující prostřednictvím místního spolku Pokrok.
Právě spolek Pokrok vydal roce 1898 rezoluci požadující rovné volební právo, ochranu chudých, zřízení burzy práce.
Městský výbor sice požadavkům nedával velkou váhu, spolky však měli i tak dopad dění ve Třebíči. \footnote{srov. \cite{Janak1981}: 97--115}

Roku 1885 je do města natažena železniční trať, průmyslová revoluce je od této chvíle v Třebíči v plném proudu.
V Třebíči ihned začali vznikat různé závody zaměřující se především na již pro Třebíč tradiční textilní a kůži zpracující průmysl, jako byl závod Budischowski.
Právě Budischowkiho zakoupil v roce 1931 Tomáš Baťa a areál rozšířil.
Počátkem 20. století je uvedena do provozu místní nemocnice, parní pivovar, sodovkárna ZON.

Během druhé světové války ve městě sídlila německá posádka.
Došlo také k deportaci 281 Židů do koncentračního tábora v Terezíně, ze kterého poté byli převáženi do vyhlazovacích táborů na východě.
Po válce se do tábora vrátilo jen pouze 10 Židů a \textit{tradice} židovského města definitivně skončila.
Zároveň se opět začali objevovat hlasy dělnictva spojené s vzrůstem komunistické strany.
Po Únoru 1948 byly zestátněny třebíčské podniky a dalšího rozmachu se městu dostane až se stavbou jaderné elektrárny Dukovany v 70. letech, v Třebíči jsou tehdy budovány nové školy a sídliště.
Po sametové revoluci je ve městě zřízena průmyslová zóna a je vyvíjena snaha opravu historických památek, právě dvě se 3. července 2003 staly památky UNESCO a to bazilika a židovské město.\footnote{srov. Historie města Třebíč v datech [online] (2006). stránka města Třebíč, Dostupné z: \url{https://www.trebic.cz/historie\%2Dmesta\%2Dv\%2Ddatech/d-1385/p1=30143} [cit. 2020-03-08]. }


\chapter{Dějiny židovského města v Třebíči}

Město Třebíč je svou židovskou čtvrtí unikátní, na rozdíl od většiny českých měst kde žily křesťané a židé více či méně odděleně, zde docházelo k těsnému soužití.
Dokonce i synagoga a bazilika jsou od sebe vzdáleny coby kamenem dohodil.
S postupem času se židovské město dostalo do obležení křesťanským Třebíčí se všech stran.
Denně se tak zde střetávaly dva značně odlišné světy a to někdy v mírů a někdy v nepřátelství.

Bohužel nám tento unikátní svět zůstává do značné míry zahalen.
Ze středověku se dochovalo jen velmi málo pramenů židovské město popisujíc.
Až úřední dokumenty z 16. století nám poskytují možnost lépe poznat osudy města.\footnote{srov. \bibentry{Fiser2005}: 10}
Situaci navíc komplikuje nejednotnost moderních prací, které se mnohdy rozchází v klíčových momentech.
Historikové se navíc vzájemně obviňují z udávání mylných dat a i dochované prameny často označují za zcestné.
Zpracovat tak rané dění židovského osídlení a vyvést absolutní závěr se tak v dnešní době jeví takřka nemožné.

\section{Založení židovské obce až utvoření ghetta}

Nejstarší datum vzniku třebíčského osídlení se udává do roku 933, ve kterém údajně byla vybudována synagoga.
Původně se však byl vznik zasazen do roku 978 dle historika J. E. Horkého, toho však J. Kořatek \uv{opravil} na již zmíněný rok 933.
Oba z těchto údajů však postrádají naprosto logiku, před založením benediktinského kláštera se na území dnešního Třebíče nacházel Třebečský les a území nebylo kolonizováno.
Právě klášter založený až roku 1101 měl sloužil jako kolonizační síla.

Někteří badatelé spojují příchod židů s tvorbou tržních osad v blízkosti kláštera.
Ani zde však nemáme žádné důkazy potvrzující vznik židovské čtvrti.
Historikové podporující tuto teorii vycházejí z důležitosti cesty z jihu na sever na které klášter ležel a blízké vesnice Týn.
Je pravda že Týn skutečně byla vesnice, kde se Židé nacházely, vznikla však mnohem později.
I o samotném významu cesty se vedou spekulace a někteří se přiklání k jejímu maximálně regionálnímu významu.
Zde hraje dost možná i stereotyp spojovat Židy s obchodem.
Osobně mi však tato teorie přijde relativně pravděpodobná a tak je zmíněna v první kapitole.

Jako další možný letopočet se nabízí rok 1238.
Na mapě publikovaném v souboru map \textit{Die markgraffschaft Mähren topograpisch} roku 1842 Gregorem Wolnym je Třebíč zakreslena se značkou Juden vor 1238 vorhanden (židé před rokem 1238).
Tento letopočet se jeví již jako relativně reálný, v té době probíhala stavba opatství, zvláště pak kostela.
Ve městě se tak nacházela řada dělníků, čímž vzrostla atraktivnost obchodovat ve Třebíči, zároveň samotné opatství potřebovalo přístupu k hotovosti.
O existenci ghetta však můžeme jen spekulovat.\footnote{srov. \cite{Fiser2005}: 21 -- 22}

Někteří odhadují vznik židovského ghetta na rok 1300, jako důkaz používají Třebíčský machzor\footnote{modlitební kniha používána Židy na Vysoké svátky jako je Roš ha-šana a Jom kipur}, který byl právě nalezen v Třebíči a obsahuje i první grafické vyobrazení Žida na území českých zemí.
Nic však nenasvědčuje tomu (ale ani nevyvrací), že rukopis vznikl právě v Třebíči.
Jestli však byl skutečně napsán v Třebíči, tak muselo ghetto vzniknout dlouho před rokem 1300 a v době napsání být jedním z významnějších na Moravě.\footnote{srov. \bibentry{Putik2015}: 24}

Prvním spolehlivým datem je až rok 1338 v tomto roce dle kronikáře opata Neplacha dochází k pronásledování židů na českém území a Třebíč je v jeho díle uvedena.
Zároveň se schoduje Martyrologium\footnote{martyrologium je list mučedníků a svatých} z norimberské židovské pamětní knihy, ta uvádí různá jihomoravská města včetně Třebíče.

Další písemný důkaz pochází z roku 1410.
Z tohoto roku pochází zápis pořízený zřejmě v Jemnici, jedná se o výslech lapka Jana Johanna ve kterém je zmíněno okradení a zadržování třebíčského Žida.
Tento zápis se dostal do jihlavské městské knihy v rámci vyměňování výslechů mezi královskými městy.
I o výkladu této poznámky se vedou spory někteří historikové hovoří až o vypálení několika domů v Podklášteří\footnote{srov. \bibentry{Klenovsky2003}: 4}, jíní jen o zajetí jednoho či více Židů kteří z Třebíče\footnote{srov. \cite{Fiser2005}: 22}.
Zároveň se ve městské knize dočteme o zastavení opatské mitry Židem.
Obecně se rok 1410 pokládá za první definitivní důkaz existence židovského města, dočteme se o něm i v úspěšném titulu Historie Židů od Tomáše Pěkného\footnote{srov. \bibentry{Pekny2001}: 419}.

O existenci židovského ghetta během husitských válek nemáme žádné prameny, a je celkem možné že čtvrť dočasně naprosto zanikla.\footnote{srov. \cite{Fiser2005}: 22 -- 23}

Další spolehlivý důkaz existence pochází až z roku 1547, jedná se o listinu zástavního držitele třebíčského kláštera, Jana z Pernštejna.
V ní je zmíněno vypovězení židů Janem Jetřichem Černohorským na základě žádosti měšťanstva.
Protože Jan z Pernštejna listinu později obnovil můžeme spekulovat, že Židé během Jana Jetřicha ve skutečnosti město neopustily.
Otázku vypovědění město dokonce konzultovalo s ostatními moravskými městy a dekret si měšťané bedlivě střežili.
Z zápisu v urbáři z roku 1556 víme, že v Podklášteří tou dobou nepřebýval jediný Žid.
Až v urbáři z roku 1573 se dočteme o šesti rodinách v Podklášteří a dvou na Stařečce.
Z urbáře se zároveň nedočteme o žádných zvláštních povinnostech Židů oproti křesťanům.
Je nutno podotknout, že židovské osídlení netvořilo v této době ghetto jak nám napovídají dvě rodiny na Stařečce a také povolení od Buriana Osovského vést židovské krámy na městském tržišti.
Právě Burian předává Židům jejich pozemky do dědičného vlastnictví.
Komunita se stále stabilně rozvíjela a v roce 1617 žilo v Podklášteří osmnáct židovských rodin.\footnote{srov. \cite{Fiser2005}: 23 -- 25}

Při převzetí panství Adamem ml. z Valdštejna dochází z neznámých důvodů k snížení počtu rodin na třináct.
Je možné, že za odchod mohla Adamova silnější katolické víra.
V urbáři z roku 1629 se zároveň dočteme, že na Židy byla uvalena povinnost platit daň z vaření medu.
Jejich \uv{rovnoprávné} postavení tak definitivně končí.
Další omezení přišla v urbáři z roku 1636, kdy začalo být pokutováno pohřbívání cizích Židy bez spravení vrchnosti, a roku 1648 přibila povinnost účasnit se roboty.
Ke konci třicetileté války pobývala v Podklášteří posádka generála Torstenssona, zda měla nějaký dopad na židovské obyvatelstvo však nevíme.\footnote{srov. \cite{Fiser2005}:25}

Po skončení třicetileté války se začíná objevovat trend přístavních domků\footnote{Domy přistavené například na zahradě} -- těch bylo postaveno do roku 1654 dvanáct.
Celkový počet rodin zároveň narostla na 36 a představovali více než polovinu v Podklášteří.
Zároveň vzniká židovský špitál.\footnote{srov. \cite{Fiser2005}: 25}

Těsně po skončení války se začíná i budovat \uv{Přední} synagoga.
Samotná výstavba probíhá mezi léty 1639--42 pravděpodobně na místě kde se nacházela dočasná dřevěná.
Ta je později roku 1757 snížena na úroveň ostatních domů a v 19. století dvakrát představena.\footnote{srov. \cite{Pekny2001}: 419}

\subsection{Hospodářská situace třebíčských Židů}

Ekonomická situace židovské populace byla ke konci 16. století srovnatelná se zbytkem měšťanů, najdeme i zmínky i zadluženosti některých židů.
V urbářích z této doby se Židé neobjevují ani jako věřitelé, až na jednu zmínku.
Zároveň nejsou zaznamenány žádné ilegální obchody, nepočítáme li výsledky inkviziční komise, té vadil prodej masa o půstu a práce v neděli.\footnote{srov. \cite{Uhlir1978}: 89, 129}

Z řemesel Židé mohli před rokem 1618 provozovat jen koželužství, pálení kořalky, rukavičkářství a mydlářství.
Během 16. století a začátkem 17. se netypicky řada Židů živila obděláváním půdy.
Od 17. století ale zažínají přecházet k řemeslům a obchodu, s tím že předchozí omezené jsou postupně odstraňována.\footnote{srov. \cite{Klenovsky2003}: 5}

\section{Vznik židovského ghetta}

Ke zformování ghetta dochází během druhé poloviny 17. století.
Během tohoto období dochází k velkému rozmachu velikosti osídlení a jen mezi roky 1654 -- 1681 je vystaveno Židy třicet nových domů.
Navíc se radikálně mění vztah Židů a křesťanů v Podklášteří.
V roce 1669 ujednala židovská komunita se správcem Františkem Augustem z Valdštejna na přeměně roboty v každoroční plat o hodnotě 100 zl.
To přispělo ke vzdálení komunity od zbytku města.
Zároveň v témže roce došlo ke stavbě Nové (Zadní) Synagogy.
Výstavba však byla v 90. letech sabotována hejtmanem a až Jan Karel z Valdštejna v roce 1705 povolil opravu a dokončení.
Mezitím neustále pokračoval trend skupování domků v Podklášteří -- v roce 1678 byl podíl majitelů domů 56 židovských ku 19 křesťanským.
Celkově však vlastnily pouze okolo jedné třetiny půdy.

1. června 1723 došlo ke vzniku ghetta.
Stává se tak díky nařízení Jana Josefa z Vladštena, které přikázalo zbývajícím křesťanům v Podklášteří výměnu domků se Židy jinde ve městě.
Od přelomu let 1723 a 1724 se začaly pořizovat samostatné židovské pozemkové knihy opět z nařízení vrchnosti.
Již v roce 1724 požádala obec o stvrzení starých \uv{svobod}.\footnote{srov. \cite{Fiser2005}: 26}

Tou dobou zároveň v českých zemích vrcholil antagonismus vůči Židům.
Během války byly židé nezbytní díky jejich finančním zásobám, ale po vestfálském míru se začal šířit názor, že židovští podnikatelé stojí v cestě křesťanským.
O tomto podnikatelském úspěchu však můžeme hovořit jen u ghetta pražského, zbytek ghett většinou skončil zadlužen.
Jedninou oporou v tomto konfilktu byla Židům dvorská a česká královská komora, ta však vídeňskému dvoru neobstála.
Karel VI. nakonec sestavil komisi s cílem omezit počet Židů ta vydala translokační reskript a tzv. familiantský zákon.
Translokační reskript z prosince 1726 přikázal Židům odstěhovat se do zvláštních ulic.
Přírůstek populace reguloval tzv. familiantský zákon -- dle něho jen židé, kteří měli během září 1727 potomstvo byly považováni za otce rodin a mohly zůstat na území.\footnote{srov. \cite{Uhlir1978}: 88 -- 95}
Právě tyto dva zákony definitivně potvrdily tvorbu ghetta.

Díky transolkačnímu reskriptu došlo k zpomalení až zastavení staveb nových domů.
Většina domů v této době byla vícepatrová, vybavena krámy a sklady pro obchodování.

\section{Židovské ghetto a město}

Během 18. a počátek 19. století potkává židovské ghetto řada pohrom.
V roce 1742 do města vtrhly pruští husaři a po židech požadovali 10 000 zl. výplaného.
I přesto, že částka byla vyplaceno, bylo několik domů vypáleno.
Dále roku 1759 vzplál dům se zemědělskými produkty, ten zahájil požár díky němuž shořelo prý 184 domů.
Takový počet domů však v Podklášteří niky nebyl, neb počet domů nikdy nepřekročil 110 a toho dosáhl až v roce 1835.
V roce 1775 postihla město povodeň a jí následovaly napoleonské války.
Jako první přes Třebíč táhnou jednotky maršála Bernotta mířící na Slavkov, tomu Židé vyplatily 5000 zl. výpalného.
Poté roku 1809 město dobudou Francouzi -- ti si vyžádají výpalné 4000 zl.
O 12 let později (1821) vypul během robotní vzpoury anti-semitisticky motivovaný požár, který poškodil okolo 100 domů.
Jemu následovala další povodeň (1830).\footnote{srov. \cite{Fiser2005}: 27}

Z organizačního hlediska patřilo obyvatelstvo do vrchnostenské jurisdikce.
V čele samotného ghetta stála rada jíž předsedal židovský rychtář, rada byla pravidelně renovována a rychtáře volili obyvatelé ghetta.
Povinná příjmení dostaly rodiny na příkaz Josefa II v roce 1799.\footnote{srov. \cite{Uhlir1978}: 44}

4.4 1849 Židé po revolučním vývoji opět dostávají svobodu pohybu a o den později možnost svobodně se ženit a vdávat.
Židům je povolen přístup do veřejných úřadů, byl také zrušen familiantský zákon, toleranční daň a hlavně ghetta.
Poslední restrikce byly odstraněny prosincovou ústavou z roku 1867. \footnote{srov. \cite{Pekny2001}: 124}
V roce 1849 tak teoreticky končí Třebíčské ghetto a můžeme hovořit o městě\footnote{Teoreticky děleno na židovskou obec a Podklášteří, obojí však bylo obýváno Židy a tak se někdy pro celek používá pojem židovské město}.
To je samostatné a disponuje vlastní zprávou podobně jako ghetto

Po zrovnoprávnění dochází ke stěhování zámožných židů do centra Třebíče.
Dělo se tak však k velké nelibosti křesťanských obyvatel a vše vyvrcholilo protižidovskými nepokoji mezi 1. až 7. dubnem  1850.

Období 40. let se také neslo v duchu sociální nejistoty drobných výrobců, ta vyvrcholila vystoupeními namířenými hlavně proti židovské čtvrti.
Během let 48-49 došlo k dočasnému zrovnoprávnění, což vedlo i k přesun zámožných židů z židovského města na náměstí.
To však mělo za následek další bouře proti židovskému obyvatelstvu, k těm došlo mezi 1. dubnem a 7. květem 1850.
Další proti židovské nepokoje v letech 1863 a 1866 byly spojeny s úpadkem soukenictví a velmi špatnou finanční situací ve městě. \footnote{srov. \cite{Janak1981}: 77}

Nutno upozornit na to, že tou dobou vrcholilo národní hnutí a židé žijící v Třebíči (nikoliv židovském městě) spolupracovali s německou menšinou.
Protože v té době bylo volební právo omezeno majetkem Židé měli nezanedbatelnou roli ve volení představitelů do městské rady Třebíče.
To na \uv{popularitě} Židů žijící v Třebíči nijak nepřidalo.\footnote{srov. \cite{Janak1981}: 93}

Velká válka a vznik Československé republiky město nijak neovlivnily.
První pokus o sjednocení židovského města a Třebíče byl 23. prosince 1924 na základě nařízení vlády ČSR.
Po odvolání bylo však 22. prosince 1926 sloučení zrušeno.
Až v září roku 1931 podala židovská obec žádost o sloučení ke kterému tentokrát došlo.\footnote{srov. \cite{Klenovsky2003}: 8}
Židé se poté postupně asimilovali do zbytku Třebíče a projevovali se jen kulturními aktivitami jako byl mužský pěvecký spolek Zion, či sportovní klub Hakoah.

Konec židovského osídlení způsobila nacistická genocida.
Ve dnech 18. a 19. května 1942 vydaly okupační úřady rozkaz k deportaci.
Ta se týkala 281 třebíčských židů a byly původně přepraveni do koncentračního tábora v Terezině a z něho redistribuováni do vyhlazovacích táborů dál na východ.\footnote{srov. Historie města Třebíč v datech [online] (2006). stránka města Třebíč, Dostupné z: \url{https://www.trebic.cz/historie\%2Dmesta\%2Dv\%2Ddatech/d-1385/p1=30143} [cit. 2020-03-08]. }
Do Třebíče se po roce 1945 vrátilo jen deset Židů, kteří se pokusily obnovit obec.
Správu nad ní však záhy převzala Židovský náboženská obec v Jihlavě a později v Brně.\footnote{srov. \cite{Klenovsky2003}: 8}
\chapter*{Závěr}
\addcontentsline{toc}{chapter}{Závěr}

Cíl práce se i přes problémy podařilo splnit.
Bohužel se mi díky epidemii CoVid--19 nepodařilo město navštívit a získat tak mé osobní fotografie.
Ty však tvoří jen malou část práce.
Během nočního přejezdu jsem ve městě zastavil a čtvrť alespoň prošel.

Jediný problém ke zpracování představovalo období dějin města od revolučního roku 1848 po současnost.
Bohužel v tomto časovém rozsahu nejsou ani publikace ohledně samotného Třebíče.
Byla sice plánována kniha Třebíč: Dějiny města III., k jejímu vytvoření snad nikdy nedošlo, či vzniklo jen pár výtisků.
Obecně nedostatek středoškolskému studentu přístupných pramenů komplikuje popsat dějiny židovské komunity v Třebíči podrobně.

I tak si myslím, že je výsledkem této práce obstojný souhrnný pohled na historii židovské obce ve Třebíči.
Sám jsem díky této práci získal hluboký náhled do dějin nejen samotného Židovského města, ale i celého Třebíče a Moravy.
Pevně doufám, že se tohoto tématu ujme schopný historik a celé dějiny zmapuje podrobně od začátku do konce.


%%% Seznam použité literatury (bibliografie)
%%%
%%% Pro vytváření bibliografie používáme bibTeX. Ten zpracovává
%%% citace v textu (např. makro \cite{...}) a vyhledává k nim literaturu
%%% v souboru literatura.bib.
%%%
%%% Příkaz \bibliographystyle určuje, jakým stylem budou citovány odkazy
%%% v textu. V závorce je název zvoleného souboru .bst. Styly plainnat
%%% a unsrt jsou standardní součástí latexových distribucí. Styl czplainnat
%%% je dodáván s touto šablonou a bibTeX ho hledá v aktuálním adresáři.

\bibliographystyle{czechiso}    %% Autor (rok) s českými spojkami
% \bibliographystyle{plainnat}    %% Autor (rok) s anglickými spojkami
%\bibliographystyle{unsrt}       %% [číslo]

\renewcommand{\bibname}{Seznam použité literatury}

%%% Vytvoření seznamu literatury. Pozor, pokud jste necitovali ani jednu
%%% položku, seznam se automaticky vynechá.

\bibliography{literatura}

%%% Kdybyste chtěli bibliografii vytvářet ručně (bez bibTeXu), lze to udělat
%%% následovně. V takovém případě se řiďte normou ISO 690 a zvyklostmi v oboru.

% \begin{thebibliography}{99}
%
% \bibitem{lamport94}
%   {\sc Lamport,} Leslie.
%   \emph{\LaTeX: A Document Preparation System}.
%   2. vydání.
%   Massachusetts: Addison Wesley, 1994.
%   ISBN 0-201-52983-1.
%
% \end{thebibliography}


\chapter*{Online zdroje}
\noindent
Historie města Třebíč v datech [online] (2006). stránka města Třebíč, Dostupné z: \url{https://www.trebic.cz/historie\%2Dmesta\%2Dv\%2Ddatech/d-1385/p1=30143} [cit. 2020-03-08].

\listoffigures
\openright
\end{document}
